\section{Ley General de Masas: atlas genealógico, lectores bidireccionales y recta dimensional}
\label{10j-sec-ley-general-masas}

La Ley General de Masas no es una tabla de valores ni una regla de ajuste por
especie. Es la realización dimensional de una genealogía construida antes de
abrir cualquier catálogo metrológico:
\begin{equation}
 \APP\longrightarrow\TRIT\longrightarrow\TPK
 \longrightarrow\widetilde\Gamma^{\rm surv}
 \longrightarrow\Gamma^{\rm obs}
 \longrightarrow L_{108}
 \longrightarrow\sigma\in\mathbb Z^5
 \longrightarrow\chi_0
 \longrightarrow\chi_{\rm full}
 \longrightarrow\widehat{\mathcal M}
 \longrightarrow\mathcal L_M .
 \label{10j-eq-cadena-masas}
\end{equation}
APP aporta las dos hojas aritméticas y su defecto orientado; TRIT conserva
signo, umbral y régimen; el TPK selecciona, transporta y actualiza el estado
mediante
\(
 U_t=\operatorname{Upd}_t\circ\operatorname{Tra}_t\circ
 \operatorname{Sel}_t
\), conservando residuo, cociente, acarreo, hoja, orientación, frontera,
ruta, memoria y supervivencia. La magnitud física aparece únicamente al final
de \eqref{10j-eq-cadena-masas}. En particular, ni una masa conocida ni una
etiqueta de especie intervienen en la selección de la ruta.

Las constantes que aparecen abajo se reciben como publicaciones internas de
la misma genealogía HMT. En concreto, \(\pi_{\rm HMT}\), \(e_{\rm HMT}\),
\(\varphi_{\rm HMT}\) y \(\alpha_{\rm HMT}\) pertenecen a la imagen nonádica
tipada generada por APP--TRIT--TPK. Su uso posterior en la rama de masa es una
composición interna descendente; no es una inyección de valores
convencionales.

\subsection{Dominio finito, rutas y multisecciones}

\begin{definition}[Atlas orientado y cocientes de fibra]
\label{10j-def-atlas}
Sea \(\mathcal C_{81}\) el atlas celular producido por el lector de familias
del estado enriquecido. Definimos
\begin{equation}
 \mathcal R_{324}:=
 \mathcal C_{81}\times\{+1,-1\}\times\{+1,-1\}.
 \label{10j-eq-rutas-324}
\end{equation}
El primer cociente identifica celdas con la misma historia no orientada y da
el conjunto \(\mathcal F_{56}\) de familias de ruta. La estratificación por
tipo de representación da trece multisecciones
\(\mathcal M_{13}\). La carta nula \(\mathcal N\) es un sumando separado del
dominio positivo; no es una decimocuarta multisección obtenida haciendo cero
un carácter.
\end{definition}

\begin{theorem}[Censo completo y no selección puntual]
\label{10j-thm-dominio}
El dominio interno de la ley satisface
\begin{equation}
 |\mathcal C_{81}|=81,
 \qquad |\mathcal F_{56}|=56,
 \qquad |\mathcal M_{13}|=13,
 \qquad |\mathcal R_{324}|=324.
 \label{10j-eq-censo-dominio}
\end{equation}
Las multiplicidades del cociente de familias son
\begin{equation}
 41\times1,\quad10\times2,\quad2\times3,\quad1\times4,\quad2\times5,
 \qquad 41+20+6+4+10=81,
 \label{10j-eq-familias-56}
\end{equation}
y los rangos de las trece multisecciones son
\begin{equation}
 \underbrace{1}_{1\ \text{ vez}},\quad
 \underbrace{2}_{2\ \text{ veces}},\quad
 \underbrace{4}_{3\ \text{ veces}},\quad
 \underbrace{6}_{2\ \text{ veces}},\quad
 \underbrace{8}_{3\ \text{ veces}},\quad
 \underbrace{12}_{1\ \text{ vez}},\quad
 \underbrace{16}_{1\ \text{ vez}},
 \label{10j-eq-rangos-13}
\end{equation}
cuya suma con multiplicidad es \(81\). La fibra de rango uno es la fibra
electrónica central. Toda fibra no central tiene rango mayor que uno y, por
tanto, su salida natural es un operador sobre la fibra; un nombre físico no
selecciona por sí solo uno de sus autovalores.
\end{theorem}

\begin{proof}
Cada celda admite exactamente las cuatro orientaciones independientes
\((h_0,v_0)\in\{\pm1\}^2\), de donde
\(|\mathcal R_{324}|=4\cdot81=324\). Las sumas de
\eqref{10j-eq-familias-56} dan simultáneamente cincuenta y seis clases y
ochenta y una celdas. La suma de rangos de \eqref{10j-eq-rangos-13} es
\[
 1+2\cdot2+3\cdot4+2\cdot6+3\cdot8+12+16=81.
\]
La inversión celular posee un único punto fijo, el centro \((5,5)\); todas
las demás órbitas forman fibras de cardinal mayor que uno. Una elección
puntual invariante en una de esas órbitas contradiría la transitividad de la
acción de inversión sobre la fibra. Por ello sólo el centro admite de forma
canónica una reducción de rango uno.
\end{proof}

\subsection{Del defecto APP al registro entero}

\begin{proposition}[Separación tipada de \(4\) y \(54\)]
\label{10j-prop-cuatro-cincuentaycuatro}
Sea \(R\) la involución que intercambia las hojas APP y
\(P_\pm=(I\pm R)/2\). El bloque central
\begin{equation}
 B_c=7I+2R=9P_++5P_-
 \label{10j-eq-bloque-app}
\end{equation}
produce el salto espectral primordial \(9-5=4\). Los números
\begin{equation}
 2,\qquad 6=51-45,\qquad
 54=9\cdot6=459-405
 \label{10j-eq-despliegue-2-6-54}
\end{equation}
son, respectivamente, acoplamiento local, defecto de compresión y defecto
integrado. En el lector electrónico, \(54=D_1(\Gamma_e)\) cuenta además un
retorno cinemático. Ninguna de esas apariciones sustituye al salto APP \(4\).
\end{proposition}

\begin{proof}
La descomposición espectral de \eqref{10j-eq-bloque-app} da autovalores
\(9\) y \(5\), luego su diferencia es \(4\). La entrada fuera de la
diagonal en la base de hojas es \(2\). Las dos compresiones globales se
evalúan en \(51\) y \(45\), y su diferencia es \(6\); integrar ese defecto
sobre nueve posiciones da \(54\). Dominio, operación y codominio cambian en
cada paso. La igualdad numérica posterior con \(D_1(\Gamma_e)\) es una
confluencia de lectores en \(\mathbb Z\), no una identidad de los objetos
que evalúan.
\end{proof}

Para una ruta \(\Gamma\in\mathcal R_{324}\), el registro CT--108 conserva
una palabra temporal \(w_\Gamma\in\mathbb F_3^{108}\), una palabra firmada
\(\tau_\Gamma\in\{-1,0,+1\}^{12}\), hoja, devanado, memoria y acarreo. El
retorno nonádico satisface
\begin{equation}
 g(K+9)=g(K),\qquad B_{K+9}\ne B_K\quad\hbox{en general},
 \label{10j-eq-retorno-memoria}
\end{equation}
porque la fase vuelve y el registro dodecafásico avanza.

\begin{definition}[Firma entera y carácter central]
\label{10j-def-firma-caracter}
El ledger de la ruta define
\begin{equation}
 \sigma(\Gamma)=
 (n_A,s_C,k_{\Delta_4},b_{120},b_\zeta)\in\mathbb Z^5.
 \label{10j-eq-firma-z5}
\end{equation}
El cierre de doce eventos determina la lectura orientada
\(s_C^{(12)}=s_C/6\); el divisor seis no es un parámetro ajustado. Sean
\begin{equation}
 \Delta_4:=
 \frac{\pi_{\rm HMT}-e_{\rm HMT}\log\pi_{\rm HMT}}{270},
 \qquad \zeta_{270}:=135\Delta_4^2,
 \label{10j-eq-delta-zeta}
\end{equation}
\begin{equation}
 A:=1000\alpha_{\rm HMT},qquad
 x:=\frac{\pi A}{180},qquad
 y:=\frac{\pi C^*}{180},qquad
 q_\pm:=e^{-x\mp y},qquad
 s_n:=q_-^n-q_+^n .
 \label{10j-eq-canales}
\end{equation}
Aquí \(C^*\) es el contraángulo ya construido por el selector regional HMT,
no por una masa. El carácter central es
\begin{equation}
 \chi_0(\sigma):=
 \exp\!\left(
 n_Ax+\frac{s_C}{6}y+k_{\Delta_4}\Delta_4
 +b_{120}s_{120}+b_\zeta\zeta_{270}
 \right).
 \label{10j-eq-caracter-central}
\end{equation}
\end{definition}

\begin{lemma}[Homomorfismo del carácter]
\label{10j-lem-caracter}
Para firmas componibles,
\(
 \chi_0(\sigma_1+\sigma_2)=\chi_0(\sigma_1)\chi_0(\sigma_2)
\). El carácter es positivo, adimensional y no contiene una unidad de masa.
\end{lemma}

\begin{proof}
El exponente de \eqref{10j-eq-caracter-central} es una forma lineal sobre
\(\mathbb Z^5\); la exponencial transforma suma en producto. Todos sus
coeficientes son coordenadas adimensionales construidas antes de abrir una
carta dimensional. La unidad aparece sólo al pasar a \(\mathcal L_M\).
\end{proof}

\subsection{Escala absoluta de acción y corrector relativo}

\begin{theorem}[Origen determinantal de la década \(-34\)]
\label{10j-thm-escala-menos34}
Para
\begin{equation}
 H_4=
 \begin{pmatrix}
 1&1&1&1\\
 1&1&-1&-1\\
 1&-1&1&-1\\
 1&-1&-1&1
 \end{pmatrix}
 \label{10j-eq-h4}
\end{equation}
se tiene
\begin{equation}
 \operatorname{SNF}(H_4)=\operatorname{diag}(1,2,2,4),
 \qquad |\operatorname{coker}H_4|=16.
 \label{10j-eq-snf-h4}
\end{equation}
El lector determinantal publica
\begin{equation}
 \Sigma_H:=\varphi_{\rm HMT}\alpha_{\rm HMT}^{16}
 =1.0462438381047565801842292083\ldots\times10^{-34},
 \label{10j-eq-sigma-h}
\end{equation}
por lo que \(\operatorname{ord}_{10}(\Sigma_H)=-34\).
\end{theorem}

\begin{proof}
Los divisores determinantales de \(H_4\) son
\(\delta_1=1\), \(\delta_2=2\), \(\delta_3=4\) y
\(\delta_4=|\det H_4|=16\). Por el teorema de Smith, los factores
invariantes son
\(
 d_1=1,
 d_2=\delta_2/\delta_1=2,
 d_3=\delta_3/\delta_2=2,
 d_4=\delta_4/\delta_3=4
\), lo que prueba \eqref{10j-eq-snf-h4}. Sustituyendo las salidas HMT ya
generadas \(\varphi_{\rm HMT}\) y \(\alpha_{\rm HMT}\) en el lector se
obtiene el valor de \eqref{10j-eq-sigma-h}; en particular
\(10^{-34}<\Sigma_H<10^{-33}\). No se ha usado una constante SI para
seleccionar el exponente dieciséis: éste es el orden del conúcleo.
\end{proof}

Las secciones anterior y posterior al retorno viven en una única recta de
acción:
\begin{equation}
 \hbar_{\rm pre}=H_5\,10^{-34}\mathcal U_S,
 \qquad
 \hbar_{\rm ret}=\eta_{\rm ret}\,10^{-34}\mathcal U_S,
 \label{10j-eq-secciones-accion}
\end{equation}
con
\begin{align}
 H_5&=1.0545718176461544696258218294330594765\ldots,\notag\\
 \eta_{\rm ret}&=1.0545718169150390611336905847242997339\ldots,\notag\\
 R_{\rm act}&:=\frac{H_5}{\eta_{\rm ret}}
 =1.000000000693281763048512336828598947\ldots .
 \label{10j-eq-ract}
\end{align}
La diferencia \(H_5-\eta_{\rm ret}\) y el cociente \(R_{\rm act}\) son
lecturas aditiva y multiplicativa de las mismas dos secciones, no dos
constantes de Planck.

\begin{proposition}[Separación de escala y refinamiento]
\label{10j-prop-escala-corrector}
Sea \(t_0>0\) la sección temporal declarada y
\begin{equation}
 \omega_0=\frac{2\pi_{\rm HMT}}{108t_0},
 \qquad \varepsilon_{\rm clk}=\hbar_{\rm ret}\omega_0,
 \qquad m_{\rm clk}=\varepsilon_{\rm clk}c_{\rm int}^{-2}.
 \label{10j-eq-reloj-masa}
\end{equation}
Entonces \(10^{-34}\mathcal U_S\) fija la sección absoluta de la recta de
masa, mientras \(R_{\rm act}^{\kappa}\) es un corrector relativo positivo.
La década \(-34\) no forma parte de \(\kappa\) y cancela exactamente en toda
razón de masas.
\end{proposition}

\begin{proof}
Por \eqref{10j-eq-secciones-accion}, ambas secciones contienen el mismo
factor \(10^{-34}\mathcal U_S\). En el cociente se obtiene
\(\hbar_{\rm pre}/\hbar_{\rm ret}=H_5/\eta_{\rm ret}=R_{\rm act}\), de modo
que la potencia decimal cancela. Si
\(m_X=m_{\rm clk}\chi_XR_{\rm act}^{\kappa_X}\) y
\(m_Y=m_{\rm clk}\chi_YR_{\rm act}^{\kappa_Y}\), entonces
\[
 \frac{m_Y}{m_X}
 =\frac{\chi_Y}{\chi_X}
  R_{\rm act}^{\kappa_Y-\kappa_X}.
\]
Introducir \(-34\) dentro de \(\kappa\), o volver a multiplicar una
coordenada ya expresada en MeV por \(10^{-34}\), contaría dos veces la misma
escala y violaría la tipificación.
\end{proof}

\subsection{Lectores bidireccionales y censo exacto}

Para \(k\in\mathbb Z/108\mathbb Z\), defínase la discrepancia cíclica
\begin{equation}
 D_\Gamma(k):=
 \sum_{t=0}^{107}
 \mathbf1_{\{w_{\Gamma,t}\ne w_{\Gamma,t+k}\}}.
 \label{10j-eq-discrepancia}
\end{equation}
El lector cronológico y el lector dodecafásico son
\begin{align}
 K_D(\Gamma)
 &:=\left(D_{27}+D_{36},D_1,\frac{D_4-D_3}{2}\right),
 \label{10j-eq-kd}\\
 K_\Omega(\Gamma)
 &:=(\Omega_{90\mid120}(\tau_\Gamma),54,P_6(\tau_\Gamma)).
 \label{10j-eq-komega}
\end{align}
El primero lee desplazamientos de \(w_\Gamma\); el segundo, autocorrelación
de ventanas de \(\tau_\Gamma\). Comparten la ruta, pero no el cociente de
información.

\begin{theorem}[Perfiles certificados de los lectores]
\label{10j-thm-censo-lectores}
Sobre \(\mathcal R_{324}\), \(K_D\) posee exactamente catorce perfiles,
\(K_\Omega\) nueve, y el par \((K_D,K_\Omega)\) cuarenta. Hay dieciocho
rutas en las que ambos triples coinciden y el único triple coincidente es
\begin{equation}
 K_D(\Gamma)=K_\Omega(\Gamma)=(80,54,6).
 \label{10j-eq-unica-coincidencia}
\end{equation}
Las multiplicidades marginales son las de las
\cref{10j-tab-kd,10j-tab-komega}.
\end{theorem}

\begin{longtable}{@{}rrr r@{}}
\caption{Catorce perfiles de \(K_D\) sobre las 324 rutas.}
\label{10j-tab-kd}\\
\toprule
\(A_D\)&\(B_D\)&\(C_D\)&multiplicidad\\
\midrule
\endfirsthead
\toprule
\(A_D\)&\(B_D\)&\(C_D\)&multiplicidad\\
\midrule
\endhead
0&24&0&18\\
40&24&2&36\\
40&54&6&18\\
40&78&27&36\\
40&84&30&36\\
60&78&25&18\\
60&84&28&18\\
80&54&6&18\\
80&54&7&18\\
80&66&2&18\\
80&66&3&36\\
80&66&4&18\\
100&66&0&18\\
100&66&2&18\\
\midrule
\multicolumn{3}{r}{total}&324\\
\bottomrule
\end{longtable}

\begin{table}[htbp]
\centering
\caption{Nueve perfiles de \(K_\Omega\) sobre las 324 rutas.}
\label{10j-tab-komega}
\begin{tabular}{@{}rrr r@{}}
\toprule
\(\Omega\)&\(B_\Omega\)&\(P_6\)&multiplicidad\\
\midrule
80&54&6&198\\
56&54&5&68\\
48&54&5&34\\
36&54&4&8\\
20&54&4&4\\
4&54&2&4\\
40&54&4&4\\
24&54&4&2\\
24&54&2&2\\
\midrule
\multicolumn{3}{r}{total}&324\\
\bottomrule
\end{tabular}
\end{table}

\begin{proof}[Demostración del \cref{10j-thm-censo-lectores}]
Se recorre el dominio finito de \eqref{10j-eq-rutas-324}. Para cada
\((c,h_0,v_0)\), la iteración de \(U_t\) durante 108 pasos determina de
forma única \(w_\Gamma\) y \(\tau_\Gamma\); luego se evalúan
\eqref{10j-eq-discrepancia}--\eqref{10j-eq-komega}. Agrupar triples iguales
produce las catorce filas de \cref{10j-tab-kd} y las nueve de
\cref{10j-tab-komega}; ambas sumas son \(324\). Agrupar pares ordenados da
cuarenta clases. Finalmente, filtrar las filas con igualdad de los dos
triples deja dieciocho rutas, todas en la única intersección
\((80,54,6)\). El algoritmo recibe sólo celda, orientación y transición
TPK; no contiene una columna de masas ni una etiqueta PDG.
\end{proof}

Las evaluaciones escalares de los lectores, cuando corresponda aplicarlas,
son
\begin{align}
 \kappa_D(\Gamma)
 &=A_D(\Gamma)-\alpha_{\rm HMT}B_D(\Gamma)
   +C_D(\Gamma)\Delta_4,
 \label{10j-eq-kappa-d}\\
 \kappa_\Omega(\Gamma)
 &=\Omega(\Gamma)-54\alpha_{\rm HMT}
   +P_6(\Gamma)\Delta_4.
 \label{10j-eq-kappa-omega}
\end{align}
La presencia de \(\alpha_{\rm HMT}\) aquí es descendente: la constante ya
ha sido publicada por HMT antes de que el lector actúe sobre una ruta.

\subsection{Electrón basal y cierre beta sin objetivo}

\begin{theorem}[Fibra electrónica y valor rector]
\label{10j-thm-electron}
El único punto fijo central determina
\begin{equation}
 \Gamma_e=(5,5,++),
 \qquad \tau_e=+++---+++---,
 \qquad
 (D_1,D_3,D_4,D_{27},D_{36})=(54,56,68,48,32).
 \label{10j-eq-ruta-electron}
\end{equation}
Por tanto,
\begin{equation}
 K_D(\Gamma_e)=K_\Omega(\Gamma_e)=(80,54,6).
 \label{10j-eq-lectores-electron}
\end{equation}
El selector basal interno \((-22,+15)\) produce
\begin{equation}
\begin{aligned}
 e_0&:=\frac{\sqrt3}{4}
 \left[
  \exp\!\left(\frac{\varphi_{\rm HMT}}{\pi_{\rm HMT}^2}\right)
  -22\alpha_{\rm HMT}^3
 \right](1+15\Delta_4)\\
 &=0.51099892248806607250987087719134943763\ldots\ {\rm MeV}.
\end{aligned}
 \label{10j-eq-electron-basal}
\end{equation}
Con
\begin{equation}
 \kappa_e:=80-54\alpha_{\rm HMT}+6\Delta_4,
 \label{10j-eq-kappa-electron}
\end{equation}
la salida beta es
\begin{equation}
 \boxed{
 m_e^\beta=e_0R_{\rm act}^{\kappa_e}
 =0.5109989506900008086082571648379233956246348677151\ldots\ {\rm MeV}.}
 \label{10j-eq-electron-beta}
\end{equation}
\end{theorem}

\begin{proof}
De \eqref{10j-eq-ruta-electron},
\[
 D_{27}+D_{36}=48+32=80,
 \qquad D_1=54,
 \qquad (D_4-D_3)/2=(68-56)/2=6,
\]
lo que prueba la primera igualdad de \eqref{10j-eq-lectores-electron}; la
autocorrelación de la palabra dodecafásica da el mismo triple. Los enteros
del selector se construyen sin masa objetivo:
\begin{equation}
 22=396/18=27-5=24-2,
 \qquad 15=270/18=\binom62=\binom64,
 \label{10j-eq-genealogia-22-15}
\end{equation}
y la matriz local de vacancias
\(
 V_e=\left(\begin{smallmatrix}21&22\\6&7\end{smallmatrix}\right)
\)
tiene determinante \(15\). Sustituir las publicaciones internas en
\eqref{10j-eq-electron-basal} da el basal indicado. La terna común
\((80,54,6)\) determina \eqref{10j-eq-kappa-electron}; finalmente,
\eqref{10j-eq-ract} da \eqref{10j-eq-electron-beta}. El valor externo se
abre sólo después de esta evaluación, por lo que no participa en la elección
de \((-22,+15)\), de la ruta ni de la potencia.
\end{proof}

El basal \(e_0\) es una etapa y no sustituye a \(m_e^\beta\). Tampoco se
confunden la firma cruda \((24,0,12,0,-12)\), la firma par CT--108
\((-12,-4,12,-6,12)\), el origen afín \(v_e=0\), la ruta \(\Gamma_e\) y
la palabra empleada por el lector de Euler: pertenecen a tipos distintos.

\subsection{Semigrupo global y convergencia de las torres}

\begin{theorem}[Semigrupo de alturas]
\label{10j-thm-semigrupo}
Las alturas de torre forman
\begin{equation}
 \mathfrak S=\langle90,120\rangle
 =\{90,120\}\cup\{30r:r\ge6\}.
 \label{10j-eq-semigrupo}
\end{equation}
Su álgebra semigrupal admite la presentación
\begin{equation}
 k[\mathfrak S\cup\{0\}]
 \simeq k[X_{90},Y_{120}]/(X_{90}^4-Y_{120}^3).
 \label{10j-eq-algebra-semigrupo}
\end{equation}
La altura \(360\) conserva dos genealogías antes del cociente: cuatro pasos
de \(90\) y tres de \(120\).
\end{theorem}

\begin{proof}
Dividir por \(30\) reduce el problema al semigrupo \(\langle3,4\rangle\).
Sus elementos positivos son \(3,4\) y todos los enteros \(r\ge6\): según el
residuo de \(r\) módulo tres, se escribe \(r=3a+4b\) con
\(b\in\{0,1,2\}\) y \(a\ge0\). La única relación primitiva entre los dos
generadores es \(4\cdot3=3\cdot4\), que, al restaurar las alturas, produce
\(X_{90}^4=Y_{120}^3\). Antes de imponer esa relación, las palabras
\(X_{90}^4\) y \(Y_{120}^3\) retienen historias distintas.
\end{proof}

Para \(h=30r\), definimos los coeficientes relativos de torre
\begin{align}
 \eta_{30r}^D(\Gamma)
 &:=D_\Gamma(9r)-D_{\Gamma_e}(9r),
 \label{10j-eq-eta-d}\\
 \eta_{30r}^\Omega(\Gamma)
 &:=A_\Gamma(r)-A_{\Gamma_e}(r),
 \qquad
 A_\Gamma(r):=
 \sum_{j=0}^{11}
 \left(\sum_{u=0}^{r-1}\tau_{\Gamma,j+u}\right)^2,
 \label{10j-eq-eta-omega}
\end{align}
donde los índices de \(\tau\) se toman módulo doce. El carácter refinado es
\begin{equation}
 \chi_{\rm full}^{\,r}(\sigma,\eta)
 :=\chi_0(\sigma)
 \exp\!\left(\sum_{h\in\mathfrak S}\eta_h^r s_h\right),
 \qquad r\in\{D,\Omega\}.
 \label{10j-eq-caracter-full}
\end{equation}

\begin{proposition}[Convergencia absoluta]
\label{10j-prop-convergencia-torres}
Las dos series de \eqref{10j-eq-caracter-full} convergen absolutamente para
toda ruta. En consecuencia, la jerarquía no requiere truncarse a las ocho
alturas usualmente exhibidas.
\end{proposition}

\begin{proof}
Por definición, \(0\le D_\Gamma(k)\le108\), luego
\(|\eta_h^D|\le108\). Como cada suma parcial de una palabra firmada tiene
módulo a lo sumo \(r\), se obtiene
\(|\eta_{30r}^\Omega|\le24r^2\). Además
\(0<q_+<q_-<1\) y
\(|s_{30r}|\le2q_-^{30r}\). Por comparación con una serie geométrica y con
\(\sum r^2q_-^{30r}\), ambas series convergen absolutamente.
\end{proof}

La torre es completa como codominio de expansión: no es un selector físico
de especie. Un algoritmo que recibe una razón para expandirla no demuestra
qué ruta produjo esa razón.

\subsection{Los trece pares de operadores de fibra}

Para cada multisección \(M\), sea
\(
 \mathcal H_M=\bigoplus_{\Gamma\in M}\mathbb C e_\Gamma
\). Definimos
\begin{equation}
 B_M^D e_\Gamma:=\kappa_D(\Gamma)e_\Gamma,
 \qquad
 B_M^\Omega e_\Gamma:=\kappa_\Omega(\Gamma)e_\Gamma.
 \label{10j-eq-operadores-bm}
\end{equation}
Las trece fibras canónicas y sus rangos son
\begin{equation}
 \begin{array}{c|ccccccccccccc}
 M&\ell_0&\ell_2&\ell_4&\nu_1&\nu_2&\nu_3&\nu_4
  &d_1&d_2&d_3&d_4&u_1&u_3\\ \hline
 \dim M&1&8&16&2&4&6&8&2&4&6&8&4&12.
 \end{array}
 \label{10j-eq-trece-fibras}
\end{equation}

\begin{theorem}[Ley operatoria sobre las multisecciones]
\label{10j-thm-ley-operatoria}
Los veintiséis operadores de \eqref{10j-eq-operadores-bm} son
autoadjuntos. Para un operador basal positivo
\(\widehat M_M^{(0)}\),
\begin{equation}
 \widehat M_M^{\beta,r}
 =R_{\rm act}^{B_M^r/2}
  \widehat M_M^{(0)}
  R_{\rm act}^{B_M^r/2},
 \qquad r\in\{D,\Omega\},
 \label{10j-eq-ley-operatoria}
\end{equation}
es positivo y covariante bajo cambio unitario de base. En la fibra
\(\ell_0\) de rango uno se reduce al escalar electrónico
\eqref{10j-eq-electron-beta}. En una fibra de rango mayor, el espectro de
\eqref{10j-eq-ley-operatoria} es la salida completa hasta que un
push-forward físico prospectivo y una escalarización estén fijados.
\end{theorem}

\begin{proof}
En la base de rutas, \(B_M^D\) y \(B_M^\Omega\) son diagonales con entradas
reales, por lo que son autoadjuntos. Como \(R_{\rm act}>0\), el cálculo
funcional define \(R_{\rm act}^{B_M^r/2}\) como operador positivo e
invertible. Para todo \(v\),
\[
 \langle v,\widehat M_M^{\beta,r}v\rangle
 =\left\langle R_{\rm act}^{B_M^r/2}v,
 \widehat M_M^{(0)}R_{\rm act}^{B_M^r/2}v\right\rangle\ge0.
\]
Si \(U\) es unitario, el cambio
\(B\mapsto UBU^*\), \(\widehat M^{(0)}\mapsto
U\widehat M^{(0)}U^*\) lleva la salida a
\(U\widehat M^{\beta,r}U^*\), y conserva espectro, traza y determinante.
Cuando \(\dim\ell_0=1\), ambos lectores valen \((80,54,6)\) y la fórmula
es escalar. Para \(\dim M>1\), escoger un autovalor sin un morfismo físico
adicional no es una operación natural sobre la fibra.
\end{proof}

La siguiente tabla exhibe las trazas vectoriales normalizadas de los triples
antes de aplicar \(\alpha_{\rm HMT}\) y \(\Delta_4\); su diversidad prueba
que no existe un único par angular que sustituya a los lectores de ruta.
\begin{longtable}{@{}c r c c@{}}
\caption{Trazas normalizadas de \(K_D\) y \(K_\Omega\) en las trece fibras.}
\label{10j-tab-trazas-13}\\
\toprule
\(M\)&\(\dim M\)&\(\overline{\mathbf K}_M^D\)&
\(\overline{\mathbf K}_M^\Omega\)\\
\midrule
\endfirsthead
\toprule
\(M\)&\(\dim M\)&\(\overline{\mathbf K}_M^D\)&
\(\overline{\mathbf K}_M^\Omega\)\\
\midrule
\endhead
\(\ell_0\)&1&\((80,54,6)\)&\((80,54,6)\)\\
\(\ell_2\)&8&\((65,249/4,17/2)\)&\((80,54,6)\)\\
\(\ell_4\)&16&\((225/4,489/8,21/2)\)&\((141/2,54,11/2)\)\\
\(\nu_1\)&2&\((40,51,29/2)\)&\((60,54,5)\)\\
\(\nu_2\)&4&\((45,63,29/2)\)&\((61,54,5)\)\\
\(\nu_3\)&6&\((170/3,55,34/3)\)&\((206/3,54,11/2)\)\\
\(\nu_4\)&8&\((105/2,237/4,95/8)\)&\((143/2,54,45/8)\)\\
\(d_1\)&2&\((40,51,29/2)\)&\((60,54,5)\)\\
\(d_2\)&4&\((45,63,57/4)\)&\((61,54,5)\)\\
\(d_3\)&6&\((170/3,55,23/2)\)&\((218/3,54,17/3)\)\\
\(d_4\)&8&\((105/2,237/4,49/4)\)&\((149/2,54,23/4)\)\\
\(u_1\)&4&\((65,66,9)\)&\((80,54,6)\)\\
\(u_3\)&12&\((205/3,119/2,113/12)\)&\((74,54,23/4)\)\\
\bottomrule
\end{longtable}

\subsection{Sector nulo y estatuto de las realizaciones físicas}

\begin{theorem}[Bifurcación nula]
\label{10j-thm-sector-nulo}
La carta fotónica \(\mathcal N_\gamma\) y la carta gluónica
\(\mathcal N_g\) se bifurcan antes del carácter positivo. En ellas
\begin{equation}
 \kappa_{\rm null}=\mathrm{N/A},
 \qquad m_\gamma=0,
 \qquad m_g=0,
 \label{10j-eq-nulidad-gauge}
\end{equation}
sin que la extensión, la energía o la representación gauge sean triviales.
La nulidad no se obtiene imponiendo \(\kappa=0\) en
\(R_{\rm act}^\kappa\).
\end{theorem}

\begin{proof}
El carácter positivo toma valores en \(\mathbb R_{>0}\), de modo que
\(R_{\rm act}^0=1\), no cero. Por tanto, ninguna potencia dentro de esa
carta puede producir la sección nula. La realización fotónica reside en la
frontera nula electromagnética y deja dos polarizaciones transversales tras
la reducción de calibre; la gluónica reside en la representación adjunta de
color de dimensión ocho. La inclusión de ambas cartas en la sección cero de
\(\mathcal L_M\) se efectúa antes de \(\chi_0\), lo que prueba
\eqref{10j-eq-nulidad-gauge} sin borrar sus representaciones.
\end{proof}

Las realizaciones físicas posteriores se ordenan por su dominio, su
representación y su modalidad de publicación; su número documental no es el
cardinal del atlas ni define la Ley General de Masas. Su estatuto exacto es:
\begin{longtable}{@{}p{.22\textwidth}p{.27\textwidth}p{.43\textwidth}@{}}
\caption{Estatuto probatorio de las realizaciones físicas posteriores.}
\label{10j-tab-realizaciones-fisicas}\\
\toprule
clase(s)&salida interna&estatuto de la coordenada de masa\\
\midrule
\endfirsthead
\toprule
clase(s)&salida interna&estatuto de la coordenada de masa\\
\midrule
\endhead
electrón&fibra \(\ell_0\), dos lectores coincidentes&escalar beta exacto,
prospectivo y sin objetivo, \eqref{10j-eq-electron-beta}\\
muón, tau&operadores de las fibras leptónicas y evaluaciones de firma
conservadas&operador exacto; las evaluaciones escalares heredadas son
condicionadas a la firma, no selectores prospectivos\\
tres neutrinos&cuatro fibras neutrínicas, límites y operador de mezcla&no se
inventan masas de sabor; la publicación primaria es operatorial\\
seis quarks&fibras \(u_j,d_j\), generación y órbita ternaria de color&operador
exacto; cualquier escalar de contraste debe declarar esquema y escala de
renormalización\\
fotón&sección nula electromagnética con dos modos transversales&nulidad exacta
\(m_\gamma=0\) en su carta\\
gluones&sección nula en la adjunta de color, con ocho generadores&nulidad exacta
\(m_g=0\) en su carta; la brecha colectiva no es masa de polo del gluón\\
\(W,Z\)&operadores gauge masivos y firmas condicionadas&la evaluación
condicionada se conserva; el join ruta--espectro físico no se sustituye por
el valor externo\\
Higgs&operador del sector escalar&escalar físico condicionado al push-forward
de representación\\
mesón, barión&gramática de composición de registros y operador de ligadura&no
son sumas desnudas de masas constituyentes\\
Fibonacci&anillo de fusión y representación de trenzado&codominio topológico;
no hay escalar universal de masa\\
Ising/Majorana&anillo de fusión, autoconjugación y término compatible&la
neutralidad no implica por sí sola Majorana ni masa\\
\(\mathbb Z_6\)&sector topológico de orden seis&operador en su codominio; no
se identifica con una altura de torre\\
virtuales&secciones de persistencia finita&clasificación por horizonte
terminal, no masa on-shell\\
\bottomrule
\end{longtable}

\subsection{Completitud del espectro interno y del inventario comparativo}

La palabra \emph{completo} se aplica aquí a dos objetos distintos y ambos se
declaran expresamente. El espacio ambiente de los dos espectros marginales es
\begin{equation}
 \mathscr A_M:=
 \coprod_{M\in\mathcal M_{13}}
 \left(
  \operatorname{Spec}\widehat M_M^{\beta,D}
  \times
  \operatorname{Spec}\widehat M_M^{\beta,\Omega}
 \right)
 \ \coprod\ \{0_\gamma,0_g\},
 \label{10j-eq-espectro-interno-completo}
\end{equation}
con multiplicidades, proyectores, rutas, esquema y escala conservados. La
salida intrínseca no es, sin embargo, todo ese producto: es el soporte de la
medida espectral conjunta, definido en
\eqref{10j-eq-espectro-conjunto-completo}, que conserva qué dos valores
proceden de una misma ruta. El inventario externo completo empleado para el
reconocimiento posterior es
\begin{equation}
 \mathscr I_{\rm ext}
 =\mathscr I_m^{471}\coprod\mathscr I_\Gamma^{384},
 \qquad |\mathscr I_{\rm ext}|=471+384=855.
 \label{10j-eq-inventario-471-384}
\end{equation}
Cada registro de \(\mathscr I_m^{471}\) distingue masa estable, masa de polo o
estimador de resonancia; cada registro de \(\mathscr I_\Gamma^{384}\)
distingue la anchura extraída de un polo o modo abierto. La correspondencia
metrológica es la relación tipada
\begin{equation}
 \mathcal C_{\rm met}
 \subset
 \mathscr A_M\times\mathscr I_{\rm ext},
 \label{10j-eq-correspondencia-metrologica-completa}
\end{equation}
que sólo se abre después de congelar descriptor, fibra, operador, proyector,
representación, esquema y escala.

\begin{theorem}[Doble completitud sin selección retrospectiva]
\label{10j-thm-doble-completitud}
La Ley General de Masas publica una salida operatorial sobre cada una de las
trece multisecciones y sobre las (324) rutas, además de las dos secciones
nulas. El inventario de \eqref{10j-eq-inventario-471-384} individualiza, sin
omisión documental, los (855) observables externos y conserva para cada uno
su sector HMT de llegada y la realización estructural requerida. Ninguna
magnitud de ese inventario selecciona una celda, una ruta, un autovalor, una
torre o un coeficiente de la ley.
\end{theorem}

\begin{proof}
El \cref{10j-thm-dominio} agota el atlas y el
\cref{10j-thm-ley-operatoria} construye los dos operadores positivos sobre
cada fibra; el \cref{10j-thm-sector-nulo} añade las dos cartas de masa cero.
Por tanto, el soporte conjunto de \eqref{10j-eq-espectro-conjunto-completo}
cubre el dominio generador íntegro y se incluye en el ambiente de
\eqref{10j-eq-espectro-interno-completo} sin identificarlo con él. El
inventario material enumera exactamente (471) registros de masa
y (384) de anchura y asigna a cada fila identidad, unidad, incertidumbre o
límite, esquema, escala, sector y realización. Permutar o sustituir las
magnitudes externas actúa únicamente sobre el segundo factor de
\eqref{10j-eq-correspondencia-metrologica-completa}; deja invariantes
\(\mathcal C_{81}\), \(\mathcal R_{324}\), los lectores, las torres y los
operadores. Esto prueba simultáneamente exhaustividad y antecedencia causal.
\end{proof}

\subsection{CKM y PMNS como transportes distintos}

En el sector quark, la incidencia N42 produce
\begin{equation}
 M_{\rm CKM}=
 \begin{pmatrix}
 2&-5&28\\
 0&7&-25/2\\
 1&-20&-2/3
 \end{pmatrix},
 \qquad \det M_{\rm CKM}=-\frac{3857}{6},
 \label{10j-eq-matriz-ckm}
\end{equation}
y transporta el operador down a la base up mediante
\begin{equation}
 B_d^{(u)}=V_{\rm CKM}B_dV_{\rm CKM}^*,
 \qquad
 R_{\rm act}^{B_d^{(u)}}
 =V_{\rm CKM}R_{\rm act}^{B_d}V_{\rm CKM}^*.
 \label{10j-eq-transporte-ckm}
\end{equation}
El transporte cambia la base, no los autovalores, y actúa después de que los
lectores hayan construido \(B_u,B_d\).
La misma matriz de incidencia actúa sobre las coordenadas angulares ya
construidas mediante
\begin{equation}
 \begin{pmatrix}\theta_{12}\\\theta_{23}\\\theta_{13}\end{pmatrix}
 =M_{\rm CKM}
 \begin{pmatrix}
  A\\ C^*/6\\ (180/\pi_{\rm HMT})\Delta_4
 \end{pmatrix},
 \qquad \delta_{\rm CP}=9A,
 \label{10j-eq-angulos-ckm}
\end{equation}
de donde
\begin{equation}
 (\theta_{12},\theta_{23},\theta_{13},\delta_{\rm CP})
 =(13.003313589509,\,2.398056157332,\,0.213977382244,
 \,65.676173123554)^\circ
 \label{10j-eq-salida-ckm}
\end{equation}
y
\begin{equation}
 J_{\rm HMT}=3.1189723326632974\times10^{-5},\qquad
 3857\delta_{\rm CP}
 =9(1528\theta_{12}+3380\theta_{23}+801\theta_{13}).
 \label{10j-eq-falsador-ckm}
\end{equation}
La última identidad es un falsador lineal exacto del transporte, no un ajuste
de sus ángulos.

Para el sector leptónico, sean \(F_\ell,F_\nu:\mathbb C^3\to\mathcal H\)
isometrías con la misma imagen \(\operatorname{im}P_3\), sea
\(b_K\in\operatorname{im}P_3\) unitario y sea
\(Q_K=|b_K\rangle\langle b_K|\leq P_3\) su proyector de rango uno. Defínase
\begin{equation}
 \Phi_L=I+(e^{i\pi/6}-1)Q_K,
 \qquad U_{\rm HMT}:=F_\ell^*\Phi_LF_\nu.
 \label{10j-eq-unitaria-leptonica}
\end{equation}

\begin{proposition}[Distinción y completitud de los transportes CKM--PMNS]
\label{10j-prop-ckm-pmns}
La matriz \(U_{\rm HMT}\) de \eqref{10j-eq-unitaria-leptonica} es unitaria,
pero la fase de rango uno no constituye la realización física PMNS de rango
completo. El transporte PMNS exacto se define por los marcos espectrales
completos, \(U_{\rm PMNS}=F_\ell^*F_\nu\). El núcleo \(G_0\) de registro
proporciona un certificado finito adicional si y sólo si se exhiben sus
entradas y se verifica \(\det G_0\ne0\); la condición no se promueve aquí a
hecho sin ese cálculo. CKM y PMNS tienen dominios, bases y falsadores distintos
y no pueden identificarse.
\end{proposition}

\begin{proof}
Como \(Q_K\leq P_3\) es un proyector ortogonal,
\([\Phi_L,P_3]=0\) y \(\Phi_L\) es unitaria. Además
\(F_\ell F_\ell^*=F_\nu F_\nu^*=P_3\); por tanto
\[
 U_{\rm HMT}^*U_{\rm HMT}
 =F_\nu^*\Phi_L^*P_3\Phi_LF_\nu=I_3.
\]
Sin embargo, una fase de rango uno conserva un estabilizador \(U(2)\). El
contraste de sus treinta y seis permutaciones no reproduce el ángulo pequeño
de mezcla física. El fracaso de la fase mínima elimina un sustituto degenerado,
pero no abre ni rebaja la realización PMNS, porque los marcos completos
producen directamente \(F_\ell^*F_\nu\). La polar de un núcleo de registro es
un certificado alternativo únicamente cuando pasa su criterio de rango. CKM,
en cambio, transporta entre las
bases up/down de operadores quark ya construidos mediante
\eqref{10j-eq-transporte-ckm}. Esta diferencia de dominios impide fundir
ambas construcciones.
\end{proof}

\subsection{Cierre espectral conjunto y estatuto exacto de los transportes CKM--PMNS}
\label{10j-sec-cierre-espectral-mezclas}

\paragraph{Espectro conjunto con correlación de ruta.}
Para cada multisección \(M\), la etapa canónica anterior a todo acoplamiento
no diagonal publica el par
\begin{equation}
 \widehat M_{M,\mathrm{rut}}^{\beta,r}
 :=\sum_{\Gamma\in M}m_r(\Gamma)
 |e_\Gamma\rangle\langle e_\Gamma|,
 \quad
 m_r(\Gamma):=m_{\rm clk}\,
 \chi_{\rm full}^{\,r}\bigl(\sigma(\Gamma),\eta^r(\Gamma)\bigr)
 R_{\rm act}^{\kappa_r(\Gamma)},
 \quad r\in\{D,\Omega\}.
 \label{10j-eq-par-diagonal-ruta}
\end{equation}
Los dos operadores son positivos y conmutan. Existe, por tanto, una única PVM
conjunta \(E_M^{D,\Omega}\) sobre \(\mathbb R_+^2\) tal que
\begin{equation}
 \widehat M_{M,\mathrm{rut}}^{\beta,D}=\int x\,dE_M^{D,\Omega}(x,y),
 \qquad
 \widehat M_{M,\mathrm{rut}}^{\beta,\Omega}=\int y\,dE_M^{D,\Omega}(x,y),
 \label{10j-eq-pvm-conjunta}
\end{equation}
\begin{equation}
 E_M^{D,\Omega}(A)e_\Gamma
 =\mathbf 1_A\!\bigl(m_D(\Gamma),m_\Omega(\Gamma)\bigr)e_\Gamma.
 \label{10j-eq-pvm-accion-ruta}
\end{equation}
La medida imagen
\begin{equation}
 \nu_M:=\frac1{|M|}\sum_{\Gamma\in M}
 \delta_{(m_D(\Gamma),m_\Omega(\Gamma))}
 \label{10j-eq-medida-imagen-conjunta}
\end{equation}
conserva la correlación entre ambos lectores y todas sus multiplicidades. En
consecuencia, el producto cartesiano escrito en
\eqref{10j-eq-espectro-interno-completo} es el espacio ambiente; el objeto
espectral intrínseco es
\begin{equation}
 \begin{aligned}
 \mathscr A_M^{\rm joint}
 &:=
 \coprod_{M\in\mathcal M_{13}}\operatorname{supp}E_M^{D,\Omega}
 \ \coprod\ \{0_\gamma,0_g\}\\
 &\hookrightarrow
 \coprod_{M\in\mathcal M_{13}}
 \bigl(\operatorname{Spec}\widehat M_M^{\beta,D}
 \times\operatorname{Spec}\widehat M_M^{\beta,\Omega}\bigr)
 \ \coprod\ \{0_\gamma,0_g\}.
 \end{aligned}
 \label{10j-eq-espectro-conjunto-completo}
\end{equation}
La inclusión no se promueve a igualdad sin probar que todas las combinaciones
marginales pertenecen a una misma ruta. Para un canal preparado \(P_{X,a}\),
la compresión \(P_{X,a}E_M^{D,\Omega}(\cdot)P_{X,a}\) es la POVM inducida.
Hay una línea de masa \(m_X\) sólo si
\begin{equation}
 E_{\mathcal M_X}(\{m_X\})P_{X,a}=P_{X,a},
 \qquad
 (\mathcal M_X-m_XI)P_{X,a}=0.
 \label{10j-eq-criterio-linea-espectral}
\end{equation}
Una compresión con varios puntos de soporte publica un multiplete o una
distribución, no una media. Si un acoplamiento posterior deja de ser diagonal,
se conserva el par ordenado y sólo se afirma una PVM conjunta después de
verificar el conmutador nulo.

\paragraph{Transporte CKM.}
Con \(s_{ij}=\sin\theta_{ij}\), \(c_{ij}=\cos\theta_{ij}\) y
\(\delta=\delta_{\rm CP}\), la salida de \eqref{10j-eq-angulos-ckm} determina
\begin{equation}
 V_{\rm CKM}=\begin{pmatrix}
 c_{12}c_{13}&s_{12}c_{13}&s_{13}e^{-i\delta}\\
 -s_{12}c_{23}-c_{12}s_{23}s_{13}e^{i\delta}&
 c_{12}c_{23}-s_{12}s_{23}s_{13}e^{i\delta}&s_{23}c_{13}\\
 s_{12}s_{23}-c_{12}c_{23}s_{13}e^{i\delta}&
 -c_{12}s_{23}-s_{12}c_{23}s_{13}e^{i\delta}&c_{23}c_{13}
 \end{pmatrix}.
 \label{10j-eq-vckm-explicita}
\end{equation}
Es unitaria, realiza el cambio de marco de
\eqref{10j-eq-transporte-ckm} y no crea los autovalores. Si
\(\kappa_i^u\) y \(\kappa_a^d\) son los espectros de los operadores
hermíticos de tipo arriba y abajo, entonces
\begin{equation}
 \left|\det[B_u,V_{\rm CKM}B_dV_{\rm CKM}^*]\right|
 =2|J_{\rm HMT}|
 \prod_{i<j}|\kappa_i^u-\kappa_j^u|
 \prod_{a<b}|\kappa_a^d-\kappa_b^d|.
 \label{10j-eq-determinante-conmutador-ckm}
\end{equation}
Éste es el falsador espectro--mezcla interno. La inferencia estadística frente
a datos externos permanece en la capa de reconocimiento y exige su covarianza.

\paragraph{Transporte PMNS y criterio de rango completo.}
La fase mínima de \eqref{10j-eq-unitaria-leptonica} satisface
\(Q_K\leq P_3\), luego \([\Phi_L,P_3]=0\); su unitariedad es exacta, aunque
permanece como control negativo de rango uno. Sean \(M_\ell\) y \(M_\nu\) los
operadores cargado y neutro ya construidos, y sean \(F_\ell,F_\nu\) marcos
espectrales ortonormales completos, con
\(F_\ell F_\ell^*=F_\nu F_\nu^*=P_3\). El transporte espectral es
\begin{equation}
 U_{\rm PMNS}:=F_\ell^*F_\nu,
 \qquad U_{\rm PMNS}^*U_{\rm PMNS}=I_3.
 \label{10j-eq-pmns-marcos-espectrales}
\end{equation}
Si ambos espectros son simples, queda determinado salvo reajustes diagonales
de fase y el orden espectral declarado. Si hay degeneraciones, la salida
exacta es la clase doble
\begin{equation}
 \prod_\lambda U(m_{\ell,\lambda})\backslash U(3)/
 \prod_\mu U(m_{\nu,\mu}),
 \label{10j-eq-pmns-clase-doble}
\end{equation}
no una matriz numérica elegida dentro de los subespacios degenerados.

El núcleo de registro conserva acarreo, hoja, orientación, vacancia y
holonomía. Con omisión de sumandos de norma nula, puede escribirse
\begin{equation}
 \begin{aligned}
 K_{\rm reg}(c,d)&:=
 \sum_{r\in\{\mathrm{car},\mathrm{hoj},\mathrm{ori},\mathrm{vac},\mathrm{hol}\}}
 \frac{\langle\Xi_r(c),\Xi_r(d)\rangle}{\sum_x\|\Xi_r(x)\|^2},\\
 (G_0)_{ab}&:=\frac1{\sqrt{|L_a||N_b|}}
 \sum_{c\in L_a}\sum_{d\in N_b}K_{\rm reg}(c,d).
 \end{aligned}
 \label{10j-eq-pmns-kernel-completo}
\end{equation}
La condición
\begin{equation}
 \det G_0\ne0
 \quad\Longleftrightarrow\quad
 \text{el registro declarado separa tres direcciones},
 \qquad
 U_{G_0}:=G_0(G_0^*G_0)^{-1/2},
 \label{10j-eq-pmns-criterio-rango}
\end{equation}
es un criterio finito exacto. No se afirma como hecho certificado mientras no
se exhiban las entradas de \(G_0\) y el cálculo de su determinante. Si
\(\det G_0=0\), debe añadirse un observable interno independiente y repetirse
el cálculo; queda prohibido forzar la invertibilidad con ángulos PMNS externos.
La PVM conjunta, su medida imagen, el transporte CKM y el teorema de marcos
PMNS son resultados internos exactos con la estructura de partida declarada;
el criterio de \(G_0\) permanece separado de esos resultados.

\subsection{No gobernanza de las tablas históricas}

\begin{proposition}[Independencia prospectiva]
\label{10j-prop-tablas-historicas}
Las antiguas ventanas de diez firmas y doce masas son testigos históricos de
comparación. No definen \(\mathcal C_{81}\), \(\mathcal F_{56}\),
\(\mathcal M_{13}\), \(\mathcal R_{324}\), los lectores, el carácter ni las
torres. Permutar, borrar o sustituir sus columnas metrológicas deja
inalteradas todas las construcciones de
\eqref{10j-eq-cadena-masas} anteriores al contraste.
\end{proposition}

\begin{proof}
Los datos de entrada de la construcción son el atlas, las cuatro
orientaciones, la transición \(U_t\), el ledger \(81\times108\), la firma y
los lectores. Ninguna de las definiciones
\eqref{10j-eq-rutas-324}, \eqref{10j-eq-firma-z5},
\eqref{10j-eq-kd}, \eqref{10j-eq-komega} o
\eqref{10j-eq-semigrupo} contiene una masa externa. Por consiguiente, una
permutación del catálogo posterior conmuta con la proyección que olvida el
contraste y actúa como la identidad sobre el constructor interno. Una tabla
abreviada no puede ser, pues, dominio ni generador de la ley.
\end{proof}

\subsection{Teorema de cierre tipado y falsadores}

\begin{theorem}[Ley General de Masas HMT--MD]
\label{10j-thm-ley-general}
La cadena \eqref{10j-eq-cadena-masas} define sobre el dominio
\(81/56/13/324\), con sector nulo separado, una ley de masa que cumple:
\begin{enumerate}
 \item construye ruta, ledger, firma, carácter y torre sin magnitudes
 externas;
 \item fija la escala absoluta mediante el conúcleo de orden dieciséis y la
 década \(-34\), separada de todo corrector relativo;
 \item publica dos lectores no equivalentes con censos
 \(14/9/40/18\) y coincidencia única \((80,54,6)\);
 \item cierra el electrón por la reducción de rango uno
 \eqref{10j-eq-electron-beta};
 \item extiende exactamente la ley como los pares de operadores
 \((B_M^D,B_M^\Omega)\) sobre las trece multisecciones;
 \item bifurca fotón y gluones hacia cartas nulas antes del carácter;
 \item publica en cada fibra no central el operador positivo, su medida
 espectral completa y sus multiplicidades; toda coordenada escalar física es
 la imagen del push-forward y de la escalarización declarados para esa
 representación;
 \item incorpora, como capa comparativa posterior, el inventario exhaustivo de
 (471) masas y (384) anchuras sin permitir que sus valores retroactúen
 sobre el atlas;
 \item mantiene distintos CKM y PMNS, construye la salida angular CKM con su
 falsador lineal y el transporte PMNS de rango completo, y abre PDG/CODATA
 sólo después de congelar la salida interna.
\end{enumerate}
No es una ley de un solo par angular: \(A,C^*\) parametrizan los canales
globales, mientras la separación de especies reside en la ruta, el ledger,
los dos lectores, las torres, la representación y el acoplamiento.
\end{theorem}

\begin{proof}
Los puntos 1 y 2 son los
\cref{10j-lem-caracter,10j-thm-escala-menos34,10j-prop-escala-corrector};
el punto 3 es el \cref{10j-thm-censo-lectores}; el punto 4 es el
\cref{10j-thm-electron}; los puntos 5 y 6 son los
\cref{10j-thm-ley-operatoria,10j-thm-sector-nulo}; el estatuto del punto 7
está tipado en \cref{10j-tab-realizaciones-fisicas}; el punto 8 es el
\cref{10j-thm-doble-completitud}; y el punto 9 es el
\cref{10j-prop-ckm-pmns}. La
\cref{10j-prop-tablas-historicas} prueba que el contraste externo no puede
retroactuar sobre esa cadena. En conjunto, las construcciones componen la
ley afirmada sin identificar escala, carácter, lector, operador y
escalarización.
\end{proof}

\begin{falsifier}[Controles negativos de la ley]
\label{10j-falsador-ley-masas}
La construcción queda refutada o mal tipada si ocurre cualquiera de los
siguientes hechos:
\begin{enumerate}
 \item una masa, CODATA, PDG o una tabla 10/12 selecciona una celda, ruta,
 firma, coeficiente o autovalor;
 \item se identifica el salto APP \(4\) con el defecto, despliegue o retorno
 \(54\);
 \item se inserta \(-34\) en \(\kappa\) o se multiplica de nuevo una
 coordenada MeV por \(10^{-34}\);
 \item los censos de los lectores dejan de ser \(14,9,40,18\), o aparece una
 coincidencia distinta de \((80,54,6)\);
 \item el basal \(e_0\) se presenta como salida electrónica final;
 \item el semigrupo \(\langle90,120\rangle\) se reduce a una lista finita o
 se usa para seleccionar retrospectivamente una especie;
 \item una fibra de rango mayor que uno recibe un escalar sin un push-forward
 físico y una escalarización prospectivos;
 \item la nulidad del fotón o del gluón se deduce de
 \(R_{\rm act}^{0}=1\), o se confunde la brecha colectiva con una masa de
 polo gluónica;
 \item se identifican CKM y PMNS o se promueve la fase leptónica de rango uno
 a mezcla física completa;
 \item falta una fila del inventario (471+384), se confunde una anchura con
 una masa o un valor externo modifica el espectro interno;
 \item se reemplaza la cadena genealógica por un único par angular.
\end{enumerate}
\end{falsifier}

La arquitectura de la ley es
\textsc{arquitectura autoral preexistente}; la recuperación conjunta del
atlas, los lectores y el electrón es \textsc{resultado recuperado}; la
presentación operatorial sobre las trece fibras es una formalización reunida;
y los censos ejecutables son certificados. Ninguna de esas categorías altera
la precedencia causal de APP--TRIT--TPK.
